\documentclass
[kulak] % [kulak|kul|kuleuven-brugge,handout|outline|]
{kulakbeamer}

\usepackage[dutch]{babel}
\usepackage[T1]{fontenc}

\title{Werken met Git en Github vanuit VSCode}
\author{Stijn Rebry}
\date{Academiejaar 2026 -- 2027}

\definecolor{GitOrange}{RGB}{240,80,50}
\definecolor{GitHubGray}{RGB}{36,41,46}
\definecolor{VSCodeBlue}{RGB}{0,122,204}

\renewcommand{\outlineTitle}{Overzicht} % enkel met optie outline

\begin{document}
	
	\begin{titleframe}
		\titlepage
	\end{titleframe}

    \begin{frame}{Inleiding}
		\begin{itemize}
            \item \textbf{Versiebeheer}
            \begin{itemize}
                \item Informatie bewaren en herstellen
                \item Samenwerken aan hetzelfde project
                \item Gestructureerd werken als team
            \end{itemize} \pause
			\item \textbf{Git} bundelt in een repository
            \begin{itemize}
                \item projectbestanden
                \item samen met geschiedenis van aanpassingen en
                \item parallelle versies van het project
            \end{itemize} 
			\item \textbf{GitHub} bewaart repositories online
            \begin{itemize}
                \item backup van actuele, vorige en parallelle versies
                \item samenwerking
            \end{itemize}
			\item \textbf{VSCode} is de grafische interface van waaruit we Git aansturen.
		\end{itemize}
	\end{frame}

	\begin{outlineframe}[Overzicht]
		\tableofcontents
	\end{outlineframe}
	
	% ======================================================================
	\section{Local repository, commits}
	% ======================================================================
	
	\begin{frame}{Lokale map maken en openen in VSCode}
		\begin{columns}[T]
			\column{0.48\textwidth}
			\begin{block}{2.1 Windows Explorer}
			\begin{itemize}
				\item Maak map \texttt{C:\textbackslash Projecten}
				\item Maak submap \texttt{gittest}
			\end{itemize}
			\end{block}
			\begin{block}{2.2 Visual Studio Code}
			\begin{itemize}
				\item Open map \texttt{C:\textbackslash Projecten\textbackslash gittest}
				\item Maak nieuw bestand \texttt{gegevens.txt}\\
				% \begin{verbatim}
                \texttt{Naam: Stijn Rebry}\\
                \texttt{Opleiding: Wiskunde}
                % \end{verbatim}
				\item Sla bestand op
			\end{itemize}
            \end{block}
			
			\column{0.48\textwidth}
			\begin{alertblock}{Gebruik lokale map, geen cloudopslag.}
			\begin{itemize}
				\item Versieconflicten
				\item Corrupte bestanden
				\item Prestatieproblemen
			\end{itemize}
		\end{alertblock}

		\begin{exampleblock}{Nuttige sneltoetsen}
			\begin{itemize}
				\item New file: \texttt{ctrl+N}
				\item Save file: \texttt{ctrl+S}
			\end{itemize}
		\end{exampleblock}
		\end{columns}
	\end{frame}

	\begin{frame}{Lokaal repository initialiseren en eerste commit}
		\begin{columns}[T]
		
			\column{0.48\textwidth}
			\begin{block}{2.3 Source control}
				\begin{itemize}
					\item Open \texttt{Source Control}
					\item Klik \texttt{Trust}
					\item Klik \texttt{Initialize Repository}
				\end{itemize}
			\end{block}

			\begin{block}<2>{2.4 Source Control, Changes}
				\begin{itemize}
					\item Schrijf een commit message
					\item Klik \texttt{Commit}
					\item Kies \texttt{Always}
				\end{itemize}
			\end{block}
	
			\column{0.48\textwidth}
			\begin{exampleblock}{Changes}
				\begin{itemize}
					\item Lijst van aanpassingen:
					\item \texttt{gegevens.txt A}
					\item[]
				\end{itemize}
			\end{exampleblock}

			\begin{exampleblock}<2>{Graph}
			\begin{itemize}
				\item Lijst van commits
				\item Repository bevat nu één versie
			\end{itemize}
		\end{exampleblock}
		\end{columns}
	\end{frame}

\begin{frame}{Repository aanpassen}
		\begin{columns}[T]
			\column{0.48\textwidth}
			\begin{block}{2.5 Editor: \texttt{gegevens.txt}}
				\begin{itemize}
				\item Voeg woonplaats toe, sla op
				\item Bekijk \texttt{Changes} en \texttt{Graph}
				\item Maak een nieuwe commit
				\end{itemize}
			\end{block}

			\begin{block}<2>{2.6 Files toevoegen/verwijderen}
				\begin{itemize}
				\item Maak nieuwe files + commit\newline
                \texttt{brol.txt}, \texttt{steden.txt}
				\item Verwijder \texttt{brol.txt}
				\item Pas \texttt{steden.txt} aan + commit
				\end{itemize}
			\end{block}
	
			\column{0.48\textwidth}
			\begin{exampleblock}<2>{Graph}
				Elke commit toont veranderingen:
				\begin{itemize}
					\item \texttt{M} -- Modified file
					\item \texttt{A} -- Added file
					\item \texttt{D} -- Deleted file
				\end{itemize}
			\end{exampleblock}

			\begin{exampleblock}<2>{Editor}
			Klik op Modified file:
				\begin{itemize}
				\item Toevoegingen (groen)
				\item Verwijderde tekst (rood)
				\end{itemize}
		\end{exampleblock}
		\end{columns}
	\end{frame}

\begin{frame}{Werkmap -- Save -- Stage -- Commit -- Repository}
			\begin{exampleblock}{Terminologie}
				\begin{itemize}
				\item Werkmap: lokale map met bestanden
				\item Stage: selectie van wijzigingen voor volgende commit
				\item Commit: opgeslagen versie van werkmap op één moment
				\item Repository: geheel van alle commits, alle aanpassingen en versies
				\end{itemize}
			\end{exampleblock}
	\end{frame}

\section{Repository publiceren op GitHub}

\begin{frame}{Repository publiceren naar GitHub}
		\begin{columns}[T]
			\column{0.48\textwidth}
			\begin{block}{3.1 Source Control}
				\begin{itemize}
					\item Publish Branch
					\item Kies private of public
					\item Confirm
				\end{itemize}
			\end{block}

			\begin{block}{3.2 GitHub}
				\begin{itemize}
					\item Klik op \texttt{Repositories}
					\item Klik op \texttt{gittest}
				\end{itemize}
			\end{block}
	
			\column{0.48\textwidth}
			\begin{exampleblock}<2>{GitHub}
                Twee repositories momenteel identiek:
				\begin{itemize}
					\item Verifieer inhoud\newline
					\texttt{gegevens.txt}, \texttt{steden.txt}
                    \item Klik op \texttt{Commits}
				\end{itemize}
			\end{exampleblock}

			\begin{exampleblock}<2>{Source Control, Graph}
            \begin{itemize}
                \item Current item: lokaal repo
				\item Cloud-icon: online repo
				\item Zelfde node: beide repo's identiek
			\end{itemize}
		\end{exampleblock}
		\end{columns}
	\end{frame}

\begin{frame}{Lokale aanpassingen pushen naar GitHub}
		\begin{columns}[T]
			\column{0.48\textwidth}
			\begin{block}{3.3 Editor: \texttt{gegevens.txt}}
				\begin{itemize}
					\item Voeg toe\newline
					\texttt{Favoriete kleur: rood}
					\item Save + Commit
				\end{itemize}
			\end{block}

			\begin{block}<2>{3.4 Source Control}
				\begin{itemize}
					\item Klik op \texttt{Sync Changes}
				\end{itemize}
			\end{block}
	
			\column{0.48\textwidth}
			\begin{alertblock}{Commit = lokaal}
				\begin{itemize}
                    \item Current $\neq$ Cloud
					\item Refresh op Github (\texttt{F5})\newline
					\texttt{gegevens.txt} ongewijzigd
				\end{itemize}
			\end{alertblock}

			\begin{exampleblock}<2>{Sync = remote}
            \begin{itemize}
                \item Current = Cloud
				\item Refresh op Github (\texttt{F5})\newline
				\texttt{gegevens.txt} up to date
			\end{itemize}
		\end{exampleblock}
		\end{columns}
	\end{frame}

    
\begin{frame}{Remote aanpassingen pullen in VSCode}
		\begin{columns}[T]
			\column{0.48\textwidth}
			\begin{block}{3.5 GitHub}
				\begin{itemize}
					\item Edit \texttt{gegvevens.txt}\newline
					\texttt{Hobby: manuals schrijven}
					\item Klik op \texttt{Commit changes}
				\end{itemize}
			\end{block}

			\begin{block}<2->{3.6 Source Control, Changes}
				\begin{itemize}
					\item Klik op \texttt{Pull}
				\end{itemize}
			\end{block}

	    	\column{0.48\textwidth}
			\begin{exampleblock}<3>{Terminologie}
            \begin{itemize}
                \item Local repo: eigen computer
                \item Remote repo: GitHub
                \item Push: wijzigingen náár GitHub sturen
                \item Pull: wijzigingen van GitHub halen
			\end{itemize}
		\end{exampleblock}
		\end{columns}
	\end{frame}

\section{Werken met branches}

\begin{frame}{Braches maken in local repository}
		\begin{columns}[T]
			\column{0.48\textwidth}
            \begin{exampleblock}{4.1 Branch bekijken}
            \begin{itemize}
                \item VSCode: links in statusbalk
                \item GitHub: pull-down-list
			\end{itemize}
		    \end{exampleblock}
			
            \begin{block}<2->{4.2 Branch maken in VSCode}
				\begin{itemize}
					\item Klik op \texttt{main}
					\item Klik op \texttt{Create new branch}
					\item Typ ``Hobbies''
				\end{itemize}
			\end{block}

	    	\column{0.48\textwidth}
			\begin{block}<2->{Branches vergelijken}
				\begin{itemize}
					\item Nieuw bestand niet in \texttt{main}
					\item \texttt{Graph}, \texttt{All history items}
					\item Inhoud werkmap
					\item Verborgen map \texttt{.git}
				\end{itemize}
			\end{block}




		\end{columns}
	\end{frame}

\end{document}